\section{Approach and Methodology}
\label{sec:approach_methodology}

\subsection{Development Process}
To ensure the rapid delivery of a Minimum Viable Product (MVP) and allow for iterative improvements, the project adopted an Agile development methodology. This approach enabled the team to prioritize core features—specifically the discovery and decision-making functionalities—while maintaining the flexibility to incorporate user feedback. The development lifecycle was divided into focused sprints, allowing for continuous testing, refinement of the user interface, and validation of the ``Visual-First'' design philosophy.

\subsection{Collaboration Tools}
Effective communication and version control were maintained using industry-standard collaboration tools. Git and GitHub were utilized for source code management, enabling the team to handle concurrent feature development, conduct code reviews, and track issues seamlessly. Documentation, such as user flows, market research, and feature matrices, were maintained in Markdown formats within the repository to ensure all team members had access to a single source of truth.

\subsection{Decision-making Rationale}
The feature selection and design architecture were driven by the goal of solving specific user pain points, as outlined in the Product Rationale:
\begin{itemize}
    \item \textbf{Focusing on Core Values (Discover \& Decide):} The priority was not to build a static directory, but a dynamic tool that directly ends user indecision. Features were evaluated based on an Impact vs. Effort Matrix, prioritizing high-impact functionalities like the Info Card and Personal Collections.
    \item \textbf{``Visual-First'' Philosophy:} Operating on the principle that ``people eat with their eyes,'' the design heavily prioritizes vibrant food imagery over text-heavy menus. This builds trust and evokes emotion at first glance.
    \item \textbf{Solving the Paradox of Choice:} By displaying only one dish at a time and employing a gamified swipe mechanism, the application intentionally limits simultaneous options to reduce cognitive load and turn decision-making into an entertaining experience.
    \item \textbf{Retention Strategy:} The ``Want to Try'' list (Personal Collections) was implemented as the primary hook. When users invest effort into saving their favorite spots, they create a valuable personal database, significantly increasing the application's Retention Rate.
\end{itemize}
